\documentclass[10pt]{article}
\usepackage[margin=1in]{geometry}
\usepackage[T1]{fontenc}
\usepackage{lmodern,microtype,amsmath,amssymb,amsthm}
\theoremstyle{definition}
\newtheorem{assumption}{Assumption}
\begin{document}
\section*{Finite-Sum Nonconvex Optimization}
We consider minimizing $F(x)=n^{-1}\sum_{i=1}^n f_i(x)$ on a real Hilbert space $E$, where $n\geq1$ and the differentiable components $f_i:E\to\mathbb{R}$ may be nonconvex. Their gradients satisfy a mean-square Lipschitz bound with constant $L>0$. At a given starting point $x^0\in E$, the objective lies within $\Delta\geq0$ of a global lower bound. The oracle returns $\nabla f_i(x)$ at a queried index $i$ and point $x$. For $\varepsilon>0$, the goal is a measurable output $\bar{x}\in E$ with $\mathbb{E}_{P}[\|\nabla F(\bar{x})\|^2]\leq\varepsilon^2$, where $P$ is the law of the algorithm's randomness.

The optimization problem is
\[
\inf_{x\in E} F(x).
\]
The norm $\|\cdot\|$ is induced by the inner product on $E$, which is equipped with a measurable structure. The algorithm receives $E$, $n$, $L$, $\Delta$, $\varepsilon$, and $x^0$.

Our assumptions:
\begin{assumption}[Mean-square smoothness]
Each component $f_i$ is Fr\'echet differentiable everywhere on $E$, and
\[
\frac{1}{n}\sum_{i=1}^n
\|\nabla f_i(x)-\nabla f_i(y)\|^2
\leq L^2\|x-y\|^2,
\qquad x,y\in E.
\]
Consequently, the objective gradient is
\[
\nabla F(x)=\frac{1}{n}\sum_{i=1}^n\nabla f_i(x).
\]
\end{assumption}

\begin{assumption}[Lower bound and initial gap]
There exists $\ell\in\mathbb{R}$ such that
\[
\begin{aligned}
F(x)&\geq\ell &&\text{for every }x\in E,\\
F(x^0)-\ell&\leq\Delta.
\end{aligned}
\]
\end{assumption}

Queries may use any index $i\in\{1,\ldots,n\}$ and any point $x\in E$; each query returns one exact component gradient. After the allotted oracle calls, the output $\bar{x}$ must be measurable, its squared gradient norm must be integrable under $P$, and it must satisfy
\[
\mathbb{E}_{P}\bigl[\|\nabla F(\bar{x})\|^2\bigr]
\leq\varepsilon^2.
\]
\end{document}
